\weeknum{4}
\topic[Gravity]{Gravity: Density, Structures, and Non-Uniqueness}

\workshopstart

\vspace{-0.75em}
\noindent\textbf{Duration:} 3 hours \hfill \textbf{Setting:} Workshop room (whiteboards; projector)

\section*{Preparation}

\begin{description}
  \item[Pre-reading:] Course Notes, Ch. 6
  \item[Lecture videos:]
  \item[Notes:]
\end{description}

% ---------------------------------------------------------
\section*{Purpose}
This workshop develops intuition about how density varies in porous and compacted materials, how different subsurface structures relate to density contrasts, and how the buried sphere solution illustrates non-uniqueness in gravity. Students work from physical reasoning, analytic solutions, and qualitative cross-section analysis. The session prepares students for Week~5 (filters, upward/downward continuation, reduction-to-the-pole).

% ---------------------------------------------------------
\section*{Learning Outcomes}
\begin{itemize}
  \item Estimate bulk density of porous materials using simple mixing laws.
  \item Explain how compaction and pore-fluid substitution change density.
  \item Predict density contrasts for various geologic structures.
  \item Interpret gravity anomaly width--depth relationships for the buried sphere.
  \item Recognise non-uniqueness and explain why lenticular shallow bodies can mimic deep spheres.
\end{itemize}

% ---------------------------------------------------------
\section*{Session Flow (\timelinelength\ min)}

Mixed groups of 4--5 with geology and physics/engineering backgrounds.  

\timeline[slot]{Warm-up --- What does gravity see?}{15}

In groups, rank a series of geological changes according to the density contrasts they are likely to produce. Consider examples such as felsic versus mafic rocks, porous versus compacted sediments, and common metamorphic transformations. Discuss which changes would be most visible to a gravity survey and which might produce similar gravity responses despite having very different geological origins.

\emph{Key idea:} Gravity does not measure rock names or geologic processes directly—it responds to density contrasts. Different geological processes can therefore produce similar gravity anomalies, introducing non-uniqueness into interpretation.

\timeline[slot]{Mini-lecture}{15}

\timeline[item]{Density and Porosity}{40}

Examine how porosity evolves during burial and how compaction changes bulk density. Using real porosity--depth data, develop a simple compaction model and explore how the resulting density changes influence gravity interpretation.

\timeline[item]{Geologic Structures to Gravity Anomalies}{40}

In this activity, you will discuss in groups some basic structures and how they connect to geology and density contrast.  You will then make an educated guess what the gravity anomaly in response to the structure will be.

\timeline[item]{Buried Sphere --- Width to Depth}{40}

Work with the buried sphere analytic solution to explore non-uniqueness.

Explain why increasing depth increases width but decreasing $\Delta\rho$ does not.  Discuss non-uniqueness: why can shallow, lens-shaped bodies (``lenticular'') mimic deep spheres?

\timeline[slot]{Wrap-up}{10}

\begin{enumerate}
  \item Which factor (porosity, fluid, composition) most strongly affects density in shallow basins?
  \item Which structural density contrast is hardest to distinguish with gravity alone?
\end{enumerate}

\timeline[slot]{Preview: Isostasy}{10}

This week we explored how density contrasts produce gravity anomalies. A natural next question is whether large topographic features such as mountain belts should therefore always produce large positive gravity anomalies. Surprisingly, they often do not. Next week we will investigate isostasy---the idea that columns of rock tend toward gravitational balance. We will compare the structure of oceanic and continental lithosphere, examine why topography requires compensating mass variations at depth, and derive the isostatic gravity correction. Along the way, we will see why crustal thickness, mantle density, and lithospheric structure are all important for understanding Earth's large-scale gravity field. Isostasy also provides our first look at how gravity data can constrain deep Earth structure without direct sampling.

% ---------------------------------------------------------
\workshopend
\notepage
\notepage

\subimport{activities/}{activity_porosity}
\subimport{activities/}{activity_structures}
\subimport{activities/}{activity_buried_sphere}